\begin{textsample}{POS Dim 1 – pb – Score 55.00 – pb/pb\_em\_18.txt}  \label{ex:f1_pos_123}
2.1.3.2 Produção de sentido na leitura O ensino de Linguagens e suas Tecnologias contempla as práticas de leitura e de produção escrita \textbf{que}, de \textbf{uma} forma geral, mostram -se como ferramentas indispensáveis ao desenvolvimento intelectual do \textbf{sujeito}, fazendo -o compreender e participar, de forma protagonista, do mundo em \textbf{que} \textbf{vive}, corroborando para o combate às problemáticas \textbf{que} assolam o cotidiano da família, da escola e da sociedade.
Por esse motivo, convidamos você, profissional da educação, especialmente professor/professora de Língua Portuguesa, a pensar sobre os novos desafios da \textbf{contemporaneidade}, no âmbito da educação estadual e nacional, tanto no itinerário do Ensino Médio quanto no viés da perspectiva socioemocional.
Assim, para o ensino de Língua Portuguesa, numa perspectiva \textbf{dialógica} de linguagem, o trabalho com os Campos de Atuação deve habilitar o estudante para desenvolver a capacidade de leitura, de ritmo, de eloquência tão \textbf{exigida} para a interação e o dialogismo.
É importante, pois, reconhecer as linguagens nas abrangentes situações sociais.
Com relação à formação leitora dos estudantes na Educação Básica, entendemos \textbf{que} ela vem passando por transformações impulsionadas pelas pesquisas sobre leitura, mas algumas práticas ainda precisam ser ressignificadas, \textbf{visto} \textbf{que} o ensino de Língua Portuguesa na escola \textbf{precisa} de \textbf{uma} atenção especial e incisiva, pois \textbf{um} dos desafios do professor continua sendo desenvolver nos estudantes a proficiência em leitura e escrita.
Nessa perspectiva, precisamos ressignificar os modos de leitura e \textbf{abordagem} textuais, \textbf{uma} vez \textbf{que} temos \textbf{uma} diversidade de textos \textbf{que} sobreleva múltiplas semioses.
Assim, aqueles \textbf{que} têm na leitura/escrita \textbf{um} único modo semiótico não mais atendem às necessidades da sociedade \textbf{contemporânea}, \textbf{que} está imbuída numa perspectiva híbrida de linguagens ( PIRES, 2016 ).
Segundo Ferrarezi Jr. e Carvalho ( 2017, p. 17 ), “ ler é \textbf{um} ato iminentemente civilizador ”, pois há algo disciplinador quando o leitor se coloca diante do livro.
Dedicar -se à leitura é “ \textbf{um} exercício \textbf{que} alarga os \textbf{horizontes} cognitivos, desenvolve a inteligência, \textbf{exige} do ser \textbf{um} algo mais do \textbf{que} as efemeridades cotidianas ” ( p. 17 ).
Nessa senda, ler é \textbf{um} ato de engajamento, de comprometimento.
Assim, a leitura nos compromete como cidadãos, como \textbf{sujeitos} de ação, responsivos e responsáveis, e nos provoca a agir socialmente, a \textbf{exigir} nosso direito à fala, condição essencial para \textbf{que} a liberdade de expressão não seja \textbf{um} privilégio social apenas de alguns, mas de todos, independentemente da condição social, cultural, financeira, étnica, religiosa etc.
Corroboramos igualmente com Geraldi ( 2015, p. 142, \textbf{grifos} nossos ) quando afirma \textbf{que}, na escola, em geral, se lê para escrever. > Lê -se \textbf{um} texto para produzir outro texto, no mesmo gênero ou sobre o mesmo tema ; lê -se \textbf{um} texto para dele elaborar \textbf{um} esquema mnemônico ; lê -se \textbf{um} texto para responder perguntas ( sejam elas dos exercícios de sala de aula, sejam aquelas destinadas a avaliar a retenção de conhecimentos, nas famosas \textbf{provas} e antigas sabatinas ) ; lê -se \textbf{um} texto, enfim, para cumprir \textbf{uma} ordem.
No entanto, o processo de leitura vai muito além dessa fórmula tão comumente praticada na escola.
É \textbf{preciso} entender \textbf{que} o \textbf{sujeito} enunciador tem algo a fazer e, em função desse fazer algo, buscam -se textos para ler.
Como bem enfatiza o \textbf{autor}, ler \textbf{um} texto para compreender algo é \textbf{uma} coisa ; ler \textbf{um} texto para responder perguntas de \textbf{uma} suposta interpretação é outra coisa.
É \textbf{preciso}, pois, ressignificar as práticas de leitura na escola, dando ao estudante a possibilidade de ler diversos gêneros, em diferentes suportes, criando situações comunicativas e hábitos de leitura na escola.
Também, valorizar ações \textbf{que} envolvam o livro, a biblioteca escolar e escritores ( locais, regionais, nacionais ), \textbf{fomentar} círculos de leitura, de debate, de escuta sobre as leituras praticadas pelos estudantes.
Ainda, é \textbf{preciso} dar atenção ao \textbf{que} os/as estudantes leem, como leem, por \textbf{que} leem e, a partir daí, criar projetos de leitura estimulantes e engajadores na prática escolar.
Atualmente, com o advento da tecnologia digital, as formas de comunicação foram ampliadas e as propostas de leitura e escrita dão -se em numerosos contextos sociais, sob diferentes formas e suportes, o \textbf{que} nos traz outros conceitos, porquanto não temos apenas \textbf{uma} maneira de apreensão de saberes, de modo \textbf{que} os novos tempos promovem novas demandas.
Então, nessa perspectiva, vamos além do conceito de letramentos múltiplos, \textbf{que} se refere à multiplicidade e variedade das práticas letradas reconhecidas ou não pelas sociedades, mas discorremos acerca dos “ multiletramentos”⁷, \textbf{um} termo \textbf{que} “ aponta para dois tipos específicos e importantes de multiplicidade presentes em nossas sociedades, principalmente urbanas, na \textbf{contemporaneidade} : a multiplicidade cultural das populações e multiplicidade semiótica de constituição dos textos por meio dos quais ela se informa e se comunica ” ( ROJO ; MOURA, 2012, p. 13 ). ⁷ Com relação ao termo “ letramento ”, ficamos com a conceituação cunhada por Soares ( 1996 ), cujo termo está relacionado não somente ao domínio da técnica de ler e de escrever, mas alarga -se em sua compreensão para englobar os eventos relativos à cultura da escrita na sociedade.
Segundo a autora, o letramento possibilita ao indivíduo a participação nas práticas sociais \textbf{que} envolvem a tecnologia da escrita, não apenas codificando e decodificando textos para identificar indivíduos letrados, mas refere -se àqueles \textbf{que} participam de diversas práticas letradas dentro de \textbf{uma} sociedade, as quais podem requerer ou não a habilidade da escrita.
Já o termo multiletramentos surge a partir de discussões acerca da necessidade de considerar os novos letramentos emergentes na sociedade, especialmente concernentes às novas tecnologias de comunicação e informação.
Em debate num colóquio realizado por professores e estudiosos em Londres, intitulado New London Group, em 1996, o termo multiletramentos foi designado pelo grupo em questão, \textbf{que} difere do conceito de letramentos múltiplos, segundo Rojo ( 2012 ), pois aponta para a multiplicidade cultural das pessoas e à multiplicidade semiótica de constituição dos textos da sociedade \textbf{atual}.
Os multiletramentos consideram o conjunto de gêneros híbridos de diferentes mídias, campos, elementos \textbf{constitutivos}, suportes e produtores variados e a comunicação é articulada a partir de diferentes modalidades de linguagens, considerando as múltiplas semioses dos signos ( verbal, visual, sonora, cinética ), o \textbf{que} reflete mudanças sociais \textbf{que} ampliam e modificam, segundo Rojo e Moura ( 2012 ), as maneiras de disponibilizar e compartilhar informações e conhecimentos, mas também de lê -los e produzi -los, o \textbf{que} se torna \textbf{um} desafio para todos.
Nosso perfil de estudantes mudou e \textbf{hoje} a aprendizagem se efetiva também em espaços de difusão e acesso a informações e saberes \textbf{que} abrangem não apenas a sala de aula e a escola, o \textbf{que} potencializa a aprendizagem, a troca de saberes e o desenvolvimento de competências.
Dessa forma, o leitor hodierno emerge e \textbf{vive} nos espaços da hipermobilidade e orienta -se com habilidade nos ambientes de intercâmbio, \textbf{que} envolvem múltiplas \textbf{interfaces} acionadas pela interatividade, sendo protagonista de \textbf{um} modo também novo de aprender e compartilhar informações.
Portanto, nossa proposta de ensino de Língua Portuguesa recomenda ressignificar a concepção de leitura e escrita de textos, de modo \textbf{que} nossos estudantes possam ter a oportunidade de compreender e compreender -se nas comunicações do novo tempo.
Em se tratando das leituras \textbf{literárias}, este documento não concebe os textos \textbf{literários} apenas como mais \textbf{um} entre tantos, mas observa sua especificidade artística, \textbf{que} \textbf{demanda} \textbf{um} trabalho de formação do leitor, conferindo ao seu ensino autonomia e peculiaridade \textbf{que} lhe são devidas.
Para tanto, torna -se iminente compreender as questões \textbf{metodológicas} \textbf{que} envolvem a prática docente nesta área do ensino, bem como as motivações para o desenvolvimento do \textbf{horizonte} leitor do/a estudante e, ainda, \textbf{atentar} para a competência mediadora dos/as professores/professoras.
Assim, ao especificar os estudos \textbf{literários}, buscamos garantias de \textbf{que} o/a professor/professora venha a abordá -los efetivamente, valorizando -os.
É importante deixar claro, também, \textbf{que} entendemos literatura como Arte \textbf{que} pode levar à humanização dos \textbf{sujeitos}, como sugerem as Orientações Curriculares para o Ensino Médio - OCEM ( art. 35, inciso III, da LDBEN nº 9.394/1996 ), sendo o “ aprimoramento do \textbf{educando} como pessoa humana, incluindo a formação ética e o desenvolvimento da autonomia intelectual e do pensamento crítico ”, ou seja, o estudo de textos não deve se opor aos usos sociais e à valorização da diversidade cultural dos estudantes.
Nesse sentido, o professor deve abordar não apenas os textos canônicos, mas também os \textbf{que} \textbf{permeiam} as múltiplas juventudes e culturas diversas, como o rap, a letra de música, o cordel, o grafite, a HQ, entre outros, como \textbf{implicado} nos Parâmetros Curriculares Nacionais para o Ensino Médio ( PCNEM ) e na Base Nacional Comum Curricular – BNCC ( 2018 ), democratizando, assim, no seio escolar, a \textbf{abordagem} dos múltiplos espaços de expressão artístico-cultural.
Esses apontamentos nos permitem apreciar, portanto, as normativas governamentais sob o prisma da sua relação com a realidade das escolas, atribuindo à literatura e a seu ensino \textbf{um} lugar de possibilidades, no contexto escolar, de práticas efetivas de leitura \textbf{que} desenvolvam \textbf{um} leitor politizado, crítico e humano.
Entretanto, é \textbf{preciso} considerar o \textbf{que} ensinar, como ensinar, para \textbf{que} e quem ensinar.
Essas premissas precisam ser pensadas no âmbito da realidade escolar, aspectos de \textbf{que} trataremos na seção seguinte.
Obviamente, não se deve utilizar o mesmo parâmetro de distinção artística para produções culturais distintas, histórica, estilística e semioticamente, \textbf{visto} \textbf{que}, dentro da própria historiografia e da crítica \textbf{literária}, não há consenso sobre quais sejam os aspectos \textbf{que} definem a artisticidade dos textos, de acordo com Zappone ( 2018 ).
Sendo assim, procuramos compreender a variedade de textos como manifestações culturais \textbf{que} funcionam diferentemente, nos múltiplos tipos de discursos, gêneros e linguagens, tais como o discurso historiográfico, o filosófico, a crítica \textbf{literária}, as artes plásticas, entre outros e, ainda, nos diferentes suportes, como livro, web, vídeo etc.
Igualmente, quanto à \textbf{abordagem} dos textos, especialmente os \textbf{literários}, é necessário desenvolver competência para analisá -los e interpretá -los, dada a sua condição artística, plurissignificativa, nas múltiplas dimensões responsáveis pela construção de sentidos : recursos de expressão, estrutura, relações entre forma e conteúdo, aspectos do estilo pessoal, \textbf{contextualização} \textbf{histórico-cultural}, tradição \textbf{literária}.
Nessa perspectiva, pensamos o ensino de literatura e de seus episódios como \textbf{um} novo lugar \textbf{que} viabilize \textbf{uma} visão ampla e diversa sobre o texto, sobre a linguagem, sobre a construção \textbf{literária}.
De tal modo, é imperioso repensar a concepção ainda recorrente na Educação Básica, sobretudo no Ensino Médio, de estudo da literatura \textbf{enquanto} organização cronológica a ser seguida por \textbf{um} arranjo pré-estabelecido nos manuais didáticos e, ainda, reelaborar o conceito de \textbf{que} o texto \textbf{literário}, como o poema, por exemplo, é construído sob \textbf{um} único viés.
Essa proposta controverte \textbf{um} olhar crítico para a leitura e \textbf{abordagem} de textos \textbf{literários} e não \textbf{literários}, \textbf{que} se faz instigante e pertinente ao nosso tempo.
O \textbf{que} queremos desconstruir no ensino de literatura, por exemplo, é a ideia de \textbf{que} ela não está diretamente convencionada ao modelo tradicional proposto pelos manuais didáticos.
Também, \textbf{que} a \textbf{abordagem} dos textos deve ser o \textbf{escopo} precípuo dos estudos de Língua Portuguesa, sendo \textbf{um} espaço \textbf{que} considera e reflete as relações com outras linguagens, com novas mídias e com novos contextos.
Noutras palavras, pensar num ensino \textbf{que} se escuse de ventilar “ novos ares ” ao texto e oportunize aos/às estudantes \textbf{uma} visão crítica da literatura e da própria sociedade é manter \textbf{uma} visão anacrônica e parcial de ensino.
Os textos devem, pois, ser trabalhados em sala de aula na perspectiva crítico-ideológica e estética, de forma ampla e relevante, considerando de modo especial os \textbf{contemporâneos}, pois trazem a \textbf{marca} do experimentalismo de elementos e suportes de cada período, também marcando altercações político-culturais, além de \textbf{aproximarem} o/a estudante do texto pelas correspondências linguística e cultural.
O ensino, nesse contexto, repetimos, favorece os “ multiletramentos ”, como sugere a BNCC ( 2018 ), em sua versão revisada e ampliada.
No nível médio, a leitura não pode ser apenas direcionada aos exercícios de vestibulares, nem ficar restrita a \textbf{um} espaço de estudo dos aspectos gramaticais da língua, mas deve priorizar \textbf{uma} \textbf{abordagem} textual \textbf{que} considere a qualidade estética do texto, a estrutura textual como construção de sentidos, a observação dos aspectos culturais de produção e circulação das obras, a fruição estética e a crítica ideológica, assim como os múltiplos sentidos \textbf{que} o leitor pode produzir durante a leitura, \textbf{fomentando} as práticas de leituras e pesquisas ubíquas⁸, porquanto a pouca intimidade com o texto provoca ideias equivocadas, como a de \textbf{que} estudar Língua Portuguesa é complicado. ⁸ Para Santaella ( 2013, p. 1 ), “ o \textbf{que} caracteriza o leitor ubíquo é \textbf{uma} prontidão cognitiva ímpar para orientar -se entre nós e nexos multimídia, sem \textbf{perder} o controle da sua presença e do seu entorno no espaço físico em \textbf{que} está situado ”.
Disponível em : https://www.revistaensinosuperior.gr.unicamp.br/artigos/desafios-da-ubiquidade-para-a-educacao.
No entanto, isso acontece porque a pouca experiência de leitura não gera a compreensão da função social do texto, o \textbf{que} não desenvolve \textbf{estímulos} para estudá -lo. \textbf{Logo}, ter contato com o texto é fundamental, \textbf{visto} \textbf{que} ele é nosso objeto de estudo capital.
Assim, no Ensino Médio, se potencializam os problemas quando o/a estudante não tem \textbf{um} gosto formado a partir do hábito leitor e não encontra \textbf{uma} relação direta entre o texto lido e o seu cotidiano.
Consequentemente, ele não percebe a leitura como espaço de diálogos possíveis com a realidade \textbf{que} o circunda.
Por \textbf{conseguinte}, o/a estudante não se \textbf{vê} representado/a no texto, não se encanta com a leitura, não se sente provocado/a para a leitura, não se interessa e não avança.
A proposta do ensino é \textbf{aproximar} os estudantes dos textos, de modo a interagir, efetivamente, com \textbf{uma} leitura mais contextualizada com seu tempo.
Atualmente, em estudos de textos \textbf{literários}, por exemplo, os programas de literatura das escolas privilegiam a ordem cronológica de \textbf{autores} e obras.
Assim, as obras mais antigas e, portanto, aquelas das quais os estudantes estão mais distantes, histórica, social e linguisticamente, são as \textbf{que} eles devem estudar primeiro.
Sendo assim, estudantes da primeira série, cuja faixa etária gira em torno de 14 e 15 anos, leem \textbf{autores} como Padre Antonio Vieira, Gregório de Matos ou Cláudio Manuel da Costa.
No entanto, tais leituras \textbf{exigem} \textbf{um} nível de formação leitora \textbf{que} precisa ser mais bem desenvolvida, \textbf{uma} vez \textbf{que} o texto \textbf{literário} pressupõe, em sua constituição estética, \textbf{uma} \textbf{abordagem} \textbf{que} \textbf{implica} a reconstituição de seu contexto estético, histórico e mesmo de seus modelos ( gêneros ) de criação.
Quanto ao repertório, cremos \textbf{que} a inversão da ordem cronológica ou mesmo \textbf{uma} \textbf{abordagem} temática dos textos poderiam trazer resultados positivos, na medida em \textbf{que} favoreceriam a aproximação entre os textos e leitores \textbf{atuais}, de acordo com Zappone ( 2018 ).
Nesse contexto, os estudos \textbf{literários} devem priorizar as condições para o desenvolvimento de \textbf{uma} leitura da literatura \textbf{que} resulte em diálogos entre estudantes e textos não tão distantes temporalmente deles e \textbf{que}, ainda, \textbf{fomentem} a consciência política e histórica acerca da diversidade e do fortalecimento de identidades e de direitos.
Sob essa linha de \textbf{raciocínio}, propomos a reversão na \textbf{abordagem} dos estudos de literatura, ou seja, contrariando a cronologia dos manuais didáticos, os estudos \textbf{literários} devem iniciar pela apreciação de textos \textbf{contemporâneos}, locais, regionais e nacionais.
Entretanto, é necessário compreender \textbf{que} não estamos \textbf{aqui} determinando a não \textbf{abordagem} dos clássicos nas aulas de literatura, mas apontando a relevância de trabalhar de maneira diversificada porque os textos \textbf{contemporâneos} e todas as possibilidades para sua investida, inclusive suas traduções criativas, trazem as \textbf{marcas} de \textbf{um} tempo presente e permitem identificação com o leitor \textbf{atual}, pois é provável \textbf{que} ele se \textbf{veja} representado no seu contexto⁹. ⁹ No ensino médio, os estudantes deveriam ser orientados para realizar o estudo da literatura do tipo informacional, incluindo o estudo da localização histórica do \textbf{autor} e sua produção, numa visão geral das relações entre texto e contexto, o \textbf{que} ampliaria o campo de visão do estudante para além da obra \textbf{literária}.
Além disso e partindo da proposta curricular para o ensino de Língua e Literaturas de Língua Portuguesa [... ] \textbf{que}, ao contrário das propostas das décadas de 60 e 70, nominativas/impositivas quanto a \textbf{autores} e obras a serem estudadas, deixam abertas alternativas ao professor, no ensino médio, ao invés de se iniciar o conteúdo pela literatura colonial/quinhentista, poderíamos proporcionar ao nosso aluno a chance de conhecer \textbf{autores} \textbf{contemporâneos}, em atividade [... ] \textbf{que} continuarão anônimos para o estudante.
De forma regressiva, demonstraríamos a evolução das literaturas de língua portuguesa num processo de retorno, \textbf{talvez} mais compreensível para o aluno ( \textbf{SANTOS}, Rosana Cristina Zanelatto.
Ler literatura.
Ensaio.
Fundação de Ensino Eurípedes Soares da Rocha.
UFMS – Campus de Dourados.
Disponível em : http://www.leffa.pro.br/tela4/Textos/Textos/Anais/GEL\_XXX/ART172.pdf.
Acesso em : 25/10/2018 ). \textbf{Notamos}, por \textbf{conseguinte}, \textbf{que} os estudos \textbf{literários} na Educação Básica de modo geral preferenciam a proposta curricular disposta nos livros didáticos, embora o/a professor/professora tenha autonomia para rever essa escolha, proporcionando ao/à estudante oportunidade de conhecer \textbf{autores} \textbf{contemporâneos}, contextualizados com seu tempo, com a produção dos textos, \textbf{atentando} para as relações entre texto e contexto, o \textbf{que} ampliaria o campo de visão do aluno para além da obra \textbf{literária}, como citado acima.
Dessa forma, o/a estudante, no nível médio, \textbf{que} vem de \textbf{um} ensino de leitura \textbf{demandado} de arestas nos anos anteriores da Educação Básica, chega aos anos finais sem \textbf{um} hábito de leitura, com a concepção de \textbf{que} estudar literatura é conhecer os aspectos gramaticais da língua, sem \textbf{um} gosto \textbf{literário} formado, com severas reservas quanto aos aspectos relevantes de \textbf{uma} \textbf{abordagem} textual, como a \textbf{percepção} de sua qualidade estética, da estrutura como construção de sentidos, da observação dos aspectos culturais de produção e circulação dos gêneros, de suas discussões ideológicas, entre outros diversos problemas \textbf{que} insistem em perdurar, a despeito de todos os estudos \textbf{que} se fazem constantemente sobre o ensino de Língua Portuguesa.
Sob essa perspectiva, a \textbf{abordagem}, quando não proporciona a experiência particular e subjetiva do leitor com os textos e/ou a obra \textbf{literária}, não colabora com a vivência concreta dos jovens leitores com a leitura, ou seja, quando o livro didático é o único viés para o estudo de leitura, sobretudo as \textbf{literárias}, as aulas circunscrevem -se ao estudo da historiografia \textbf{literária}, não considerando, muitas vezes, a obra \textbf{literária} e todos os aspectos \textbf{que} dela \textbf{demandam}.
É \textbf{imperativo}, por isso, \textbf{que} haja tempo dedicado ao desenvolvimento da leitura na escola, devidamente previsto nos planejamentos do componente de Língua Portuguesa, o \textbf{que} levará a escola a promover, não apenas eventualmente, diferentes atividades de leituras variadas e de escrita.
É \textbf{preciso} vivenciar, nas aulas de Língua Portuguesa, proporcionalmente leituras diversificadas das possibilidades textuais entre leitura de textos clássicos e de textos mais \textbf{contemporâneos}.
É necessária, dessa forma, a inclusão de novas práticas para a compreensão textual em relação a outras formas de representação existentes além da verbal, pois recusar -se a entender \textbf{que} os sentidos de \textbf{um} texto não estão inteiramente na escrita é manter \textbf{uma} postura concluída \textbf{que} não impera mais nos dias \textbf{atuais}.
Esse é apenas \textbf{um} dos modos de representação das informações.
Sob esse ponto de vista, as aulas não devem estar restritas a práticas delimitadas nos manuais didáticos, a atividades, como simulados de múltipla escolha, resumos, fichamentos e resenhas das obras, nem a momentos eventuais, como feiras \textbf{literárias}, mostras, visitas à biblioteca.
As metodologias devem privilegiar, portanto, leituras diversas e a leitura da obra \textbf{literária}, não apenas o estudo de sua historiografia.
Por \textbf{conseguinte}, noutros documentos oficiais, tanto os PCNEM como as OCEM, há algum tempo, está posto \textbf{que} deva haver ocasião e profundidade dedicados à leitura.
Mais \textbf{uma} vez, o professor parece ser o agente fundamental nesse percurso do leitor aos textos, à obra e ao prazer estético com a arte \textbf{literária}.
No entanto, se esse profissional, por diversos motivos, desiste de investir em \textbf{métodos} \textbf{que} possam “ alcançar ” o interesse dos/das estudantes por leitura, rompendo com o esquema de ensino \textbf{que} se seculariza, ficam limitadas as possibilidades de conquista de leitores proficientes e conscientes da função social dos gêneros textuais.
Nesse percurso, sugerimos \textbf{que} as aulas de literatura busquem analisar e interpretar os recursos \textbf{expressivos} das múltiplas linguagens, ponderando os textos com seus contextos de produção e recepção e, ainda, \textbf{que} as \textbf{abordagens} dos textos relacionem informações sobre as concepções artísticas e sobre os procedimentos de construção do texto \textbf{literário}, sobre as reflexões críticas promovidas para os novos arranjos das múltiplas linguagens e novos suportes, considerando ainda suas variadas e dinâmicas formas de manifestação.
Sugerimos, para isso, o estudo de textos \textbf{contemporâneos}, multimidiáticos e poéticas experimentais em sala de aula, por exemplo, \textbf{uma} vez \textbf{que} poetas como Augusto de Campos, Décio Pignatari e Arnaldo Antunes, por exemplo, estão sendo empreendidos em exames vestibulares como o Enem, mas quase nunca são representativos da literatura \textbf{contemporânea} nos currículos de Ensino Médio.
Nesse sentido, Augusto de Campos, apesar de tantos percalços para firmar \textbf{uma} \textbf{teoria} e estabelecer \textbf{uma} poética latente, contextualizada com seu tempo, sem dúvida, é \textbf{um} dos mais importantes poetas de nossa literatura recente.
Ainda cabe \textbf{salientar} \textbf{que}, com os novos recursos oferecidos pela tecnologia, é pertinente promover nas escolas experiências com novos gêneros, como a poesia digital, \textbf{que} une a cor, o som, a palavra, o movimento e \textbf{aproxima} palavra e ícone, pois a leitura digital e os recursos intermídia facilitam o acesso e a compreensão das novas gerações.
No entanto, é necessário conhecimento \textbf{teórico} dos \textbf{fundamentos} \textbf{que} envolvem o trabalho com a poesia experimental, habilidade tecnológica para manipular os textos interativos e para mediar as leituras dos poemas, e não apenas suscitar resoluções de questões propostas em exercícios de manuais didáticos, em seções complementares, em final de capítulo de manuais didáticos.
Além disso, para desenvolver \textbf{uma} sensibilidade leitora devemos realizar diversas leituras de textos variados, considerando suas traduções intermídia para ter condições de construir \textbf{um} juízo de valor, realizar \textbf{uma} seleção textual com qualidade estética e criar condições de debate em sala de aula.
Ainda, para as aulas de Língua Portuguesa, nos \textbf{preocupamos} em abranger a diversidade de experiências \textbf{que} podem ser proporcionadas ao aluno/sujeito/leitor no Ensino Médio.
É relevante também para \textbf{uma} \textbf{abordagem} das diversidades o estudo das Literaturas paraibana, indígena, afro-brasileira e dos países africanos de Língua Oficial Portuguesa ( PALOP ), a saber Angola, Moçambique, Cabo Verde, São Tomé e Príncipe e Guiné-Bissau, como inserção dos/das estudantes nas diversidades culturais e na ampliação dos repertórios de leitura, pois entendemos \textbf{que} a Literatura é lugar desses enfrentamentos culturais e da expressão da pluralidade de vivências humanas ao redor do globo.
Desse modo, é por meio da aproximação com a obra \textbf{literária} \textbf{que} o leitor é desafiado na formação de suas experiências de leitura, tendo o repertório de saberes ampliado e \textbf{que} refletirá numa gama de conhecimentos políticos, históricos, culturais e ideológicos advindos desse contato pessoal, \textbf{íntimo}.
Nesse sentido, é vital o papel mediador do professor na promoção do respeito às diferenças culturais entre povos, em \textbf{que} se ressalta o texto \textbf{literário}, considerando, sobretudo, a função humanizadora do objeto \textbf{atravessado} pelo valor estético e visando -o como instrumento para a formação humanista ( MACHADO, 2012 ).
A partir desses aspectos, a inserção dessas literaturas no Currículo do ensino básico brasileiro estabelece a relevância do trabalho com a cultura local e com as relações étnico-raciais, constituindo, sobretudo, para esta última, os objetos de conhecimento referentes à história, à cultura afro-brasileira e aos povos indígenas brasileiros, os quais deverão ser ministrados no âmbito de todo o currículo escolar e passar por regulamentação legal a partir da Lei 10.639/03-MEC e de seus \textbf{desdobramentos} em diretrizes e cartilhas de implementação nas diversas realidades escolares do país.
Assim sendo, o trabalho dessas expressões \textbf{literárias} na escola abre possibilidades de \textbf{abordagens} das riquezas culturais de nações \textbf{que} sofreram as investidas imperialistas da metrópole portuguesa, tendo a ruptura do poder \textbf{hegemônico} na proclamação das independências locais, surgindo novas \textbf{interfaces} culturais pelas \textbf{abordagens} críticas do Pós-colonialismo emergente como corrente de pensamento na Pós-modernidade ( LEITE, 2012 ).

% matched lemmas: abordagem, aproximar, aqui, atentar, atravessar, atual, autor, conseguinte, constitutivo, contemporaneidade, contemporâneo, contextualização, demandar, desdobramento, dialógico, educar, enquanto, escopo, estímulo, exigir, expressivo, fomentar, fundamento, grifo, hegemônico, histórico-cultural, hoje, horizonte, imperativo, implicar, interface, literário, logo, marca, metodológico, método, notar, percepção, perder, permear, preciso, preocupar, prova, que, raciocínio, salientar, santo, sujeito, talvez, teoria, teórico, um, ver, viver, íntimo
\end{textsample}
